\section{厚生年金}

\begin{fillblank}[年金1 2022 問題1(5) ア 改題]
  A さん（40 歳）は、会社勤めの夫の扶養の範囲で今の会社で働いていたが、2022 年（令和4
年）10 月施行の短時間労働者に対する厚生年金保険等の適用拡大（※）に伴い
厚生年金保険に加入した。A さんが2022 年（令和4 年）10 月から起算して10 年間今の会社
で働いた後に退職した場合の今後の保険料および年金額を整理すると、以下のとおりとな
る。
\begin{itemize}
  \item 今の会社で働いている期間における厚生年金保険料（本人負担）は、 （　　A　　） 円（月
額）となる。
  \item 年金額については、老齢基礎年金 （　　B　　） 円（月額）に加え、老齢厚生年金 （　　C　　）
円（月額）が受け取れるようになる。
\end{itemize}



なお、解答にあたっては以下を前提とすること。
\begin{itemize}
  \item A さんは20 歳から国民年金に加入（結婚後は第3 号被保険者として加入）しており、退職
後も60 歳になるまで引き続き加入するものとする。
  \item A さんの2022 年（令和4 年）10 月以降の標準報酬月額は98,000 円で賞与はないものと
し、以降退職まで不変とする。
  \item 簡単のため、厚生年金保険法第43 条に定める平均標準報酬額に関する再評価率は
「1.000」とする。
  \item 保険料率および年金額は \warekiyear 年度のものを使用するものとし、百円未満は四
捨五入することとする。
\end{itemize}

（※）厚生年金保険の企業規模要件は以下の通り段階的に撤廃される。

\begin{table}[htbp]
    \centering
    \begin{tabular}{|c|c|} \hline 
      企業規模（常勤の従業員数で判断） &  実施時期 \\ \hline
      500人超 & 2016年10月 \\ \hline
      100人超 & 2022年10月 \\ \hline
      50人超 & 2024年10月 \\ \hline
      （　　D　　）人超 & 2027年10月 \\ \hline
      （　　E　　）人超 & 2029年10月 \\ \hline
      （　　F　　）人超 & 2032年10月 \\ \hline
      （　　F　　）人以下 & 2035年10月 \\ \hline
    \end{tabular}
  \end{table}


\end{fillblank}

%毎年更新が必要
\begin{answerbox}
  \begin{enumerate}[A: ]
    \item $98,000円 \times 0.183 \times 1/2 = 8,967円$なので百円未満を四捨五入して、9,000円。
    \item 老齢基礎年金の満額は令和8年度だと月額$70,608円$なので百円未満を四捨五入して、70,600円。
    \item $98,000円 \times 5.481/1,000 \times 10年 = 5,371.38円$なので百円未満を四捨五入して、5,400円。
    \item 35
    \item 20
    \item 10
  \end{enumerate}
\end{answerbox}

\begin{explanationbox}
  後半の類題：年金1 2021 問題1(5)\textcircled{3}
\end{explanationbox}

\begin{shortanswer}[自作 2026/11/1]
  被用者保険の適用者拡大に関して、2029年10月から適用予定の個人事業所の適用拡大に関する関する内容を、
  当分の間の経過措置も含めて簡潔に記入しなさい。
\end{shortanswer}

\begin{answerbox}
  常時5人以上の者を使用する個人事業所のうち法律で定める17業種以外の業種についても、
  2029年10月以降は被用者保険の適用事業所とする。
ただし、2029年10月時点で既に存在している事業所は当分の間、対象外とする。
\end{answerbox}

\begin{explanationbox}
  年金制度改正法（令和７年６月１３日成立）

  \url{https://www.mhlw.go.jp/stf/seisakunitsuite/bunya/0000147284_00017.html}

  \;

  類題：年金1 2023 問題2(3)ア
\end{explanationbox}

\begin{fillblank}[年金1 2021 問題1(5)\textcircled{1}]
  厚生年金保険法では、標準報酬月額の改定について次のように定められている。

  \begin{lawstructure}
    \begin{lawarticle}{第二十三条}{改定}
      実施機関は、被保険者が現に使用される事業所において継続した三月間
      （各月とも、報酬支払の基礎となつた日数が、（　　A　　）日以上でなければならない。）
      に受けた報酬の総額を三で除して得た額が、その者の標準報酬月額の基礎となつた報酬月額に比べて、
      著しく高低を生じた場合において、必要があると認めるときは、その額を報酬月額として、
      その著しく高低を生じた月の翌月から、標準報酬月額を改定することができる。

      \lawparagraph{2}
      前項の規定によつて改定された標準報酬月額は、その年の（　　B　　）月
      （（　　※　　）月から（　　※　　）月までのいずれかの月から改定されたものについては、翌年の（　　B　　）月）
      までの各月の標準報酬月額とする。
    \end{lawarticle}

    （注）（　　※　　）は出題の関係上省略している。
  \end{lawstructure}
\end{fillblank}

\begin{answerbox}
  \begin{enumerate}[A: ]
    \item 17
    \item 8
  \end{enumerate}
\end{answerbox}

\begin{fillblank}[年金1 2022 問題1(5) イ]
  老齢厚生年金の繰下げ待機中の者が厚生年金保険に加入中の場合は、在職老齢年金で支給停
止された部分は繰下げ増額の対象とならず、繰下げによる増額分は以下のとおりとなる。

\begin{flalign*}
  （ 　繰下げ対象額 \alpha 　+　 （　　\text{A}　　） 　）　\times 　増額率 \beta
\end{flalign*}

  \begin{table}[htbp]
    \centering
    \begin{tabular}{|l|l|} \hline 
      繰下げ対象額$\alpha$ & 受給権発生時点での老齢厚生年金の額$\times$平均支給率$\gamma$ \\ \hline
      増額率$\beta$ & （　　B　　）％$\times$繰下げ月数（（　　C　　）月を上限とする）\\ \hline
      平均支給率$\gamma$ & 受給権が発生した月の翌月から繰下げ請求をした月までの各月の支給率$\delta$の平均 \\ \hline
      支給率$\delta$ & $1 - (各月の在職支給停止額 \div 受給権発生時点での老齢厚生年金の額)$ \\ \hline
    \end{tabular}
  \end{table}
\end{fillblank}

\begin{answerbox}
  \begin{enumerate}[A: ]
    \item 経過的加算額
    \item 0.7
    \item 120
  \end{enumerate}
\end{answerbox}


\begin{fillblank}[年金1 2020 問題1(5)\textcircled{1}]
  臨時に使用される者（船舶所有者に使用される船員を除く）であって、日々雇い入れられる
者は、厚生年金保険の被保険者としない（ただし、（　　A　　）月を超え、引き続き使用される
に至った場合を除く）。また、臨時的事業の事業所に使用される者も、厚生年金保険の被保険
者としない（ただし、（　　B　　）月を超えて使用されるべき場合は、この限りでない）。
\end{fillblank}

\begin{answerbox}
  \begin{enumerate}[A: ]
    \item 1
    \item 継続して6
  \end{enumerate}
\end{answerbox}

\begin{fillblank}[年金1 2020 問題1(5)\textcircled{2} 改題]
  毎年（　　A　　）における全被保険者の標準報酬月額を平均した額の（　　B　　）に相当する額
が標準報酬月額等級の最高等級の標準報酬月額を超える場合において、その状態が継続すると
認められるときは、その年の9 月1 日から、健康保険法第 40 条第 1 項に規定する標準報酬月
額の等級区分を参酌して、政令で、当該最高等級の上に更に等級を加える標準報酬月額の等級
区分の改定を行うことができる旨、厚生年金保険法に規定されている。

標準報酬月額の上限は以下の通り段階的に引き上げられることになっている。
\begin{table}[htbp]
    \centering
    \begin{tabular}{|c|c|} \hline 
      標準報酬月額の上限 &  実施時期 \\ \hline
      65万円（32等級） & 現行 \\ \hline
      （　　C　　）万円 & 2027年9月 \\ \hline
      （　　D　　）万円 & 2028年9月 \\ \hline
      （　　E　　）万円 & 2029年9月 \\ \hline
    \end{tabular}
  \end{table}
\end{fillblank}

\begin{answerbox}
  \begin{enumerate}[A: ]
    \item 3月31日
    \item 100分の200
    \item 68
    \item 71
    \item 75
  \end{enumerate}
\end{answerbox}


\begin{fillblank}[年金1 2020 問題1(5)\textcircled{4}]
厚生年金保険の被保険者である夫の死亡によって遺族厚生年金の受給権を取得した当時40
歳以上65 歳未満で、遺族基礎年金が支給されない（されなくなった）妻等が65 歳未満である
ときは、遺族厚生年金の額に加算額があり、\warekiyear 年度における当該加算額は
（　　A　　）円である。
\end{fillblank}

%毎年変える
\begin{answerbox}
  \begin{enumerate}[A: ]
    \item 国民年金の令和8年度の月額$70,608円 \times 12ヶ月 \times 3/4 = 635,472円$
  \end{enumerate}
\end{answerbox}

\begin{explanationbox}
  中高齢寡婦加算のこと。

  $遺族基礎年金の額（780,900円　2004年水準） \times 改定率 \times 3/4$を加算する。
\end{explanationbox}

\begin{truefalse}[年金1 2022 問題1(4)エ]
下線部の正誤を判定し、誤っている場合は正しい内容に改めなさい。
\\

  2022 年（令和4 年）5 月より、60 歳以上65 歳未満の厚生年金保険の被保険者は個人型年金
に加入できるようになった。ただし、\underline{個人型年金または企業型年金の老齢給付金の受給権を
有する者又はその受給権を有する者であった}\\
\underline{者、特別支給の老齢厚生年金の受給権を有して
いる者}などについては、個人型年金に加入できない。
\end{truefalse}

\begin{answerbox}
\textbf{誤}


  個人型年金の老齢給付金の受給権を有する者またはその受給権を有する者であった者、特別支給の老齢厚生年金の支給繰上げを請求した者
\end{answerbox}

\begin{fillblank}[年金1 2019 問題1(5)\textcircled{1}]
  2018 年9 月以降の第1 号厚生年金被保険者、第2 号厚生年金被保険者、第3 号厚生年金被
保険者の保険料率は（　　A　　）％であり、第4 号厚生年金被保険者については2027 年4 月以
降、第1 号厚生年金被保険者と同一の保険料率に統一される。
\end{fillblank}

\begin{answerbox}
  \begin{enumerate}[A: ]
    \item 18.3
  \end{enumerate}
\end{answerbox}



\begin{fillblank}[年金1 2019 問題1(5)\textcircled{2}]
  厚生年金保険法では、標準報酬月額の定時決定ついて次のように定められている。

  \begin{lawstructure}
    \begin{lawarticle}{第二十一条}{定時決定}
      実施機関は、被保険者が毎年（　　A　　）月一日現に使用される事業所において同日前（　　B　　）月間
      （その事業所で継続して使用された期間に限るものとし、かつ、
      報酬支払の基礎となつた日数が十七日（厚生労働省令で定める者にあつては、十一日。第二十三条第一項、
      第二十三条の二第一項及び第二十三条の三第一項において同じ。）未満である月があるときは、その月を除く。）
      に受けた報酬の総額をその期間の月数で除して得た額を報酬月額として、標準報酬月額を決定する。

      \lawparagraph{2}
      前項の規定によつて決定された標準報酬月額は、その年の（　　C　　）月から翌年の（　　D　　）月
      までの各月の標準報酬月額とする。

      （以下略）
    \end{lawarticle}
  \end{lawstructure}

\end{fillblank}

\begin{answerbox}
  \begin{enumerate}[A: ]
    \item 七
    \item 三
    \item 九
    \item 八
  \end{enumerate}
\end{answerbox}

\begin{shortanswer}[年金1 2020 問題2(3)]
  公的年金に関し、 「年金制度の機能強化のための国民年金法等の一部を改正する法律」（令和2
年 6 月 5 日公布）により令和4 年4 月1 日（以下「施行日」 ）施行とされた、厚生年金の在職中
の年金受給の在り方の見直しについて、次の(1)および(2)の設問に答えなさい。

\begin{enumerate}[(1)]
  \item 新たに導入されることとなった「在職定時改定」の内容を簡記しなさい。
  \item 65 歳未満の在職老齢年金制度（低在老）に関し、見直しの内容を簡記したうえで、以下に
記載のＣさんの施行日の属する月の老齢厚生年金の月額への影響を示しなさい。
\end{enumerate}


＜Ｃさんの情報（記載はすべて施行日時点のものとする）＞
\begin{itemize}
  \item 62 歳
  \item 総報酬月額相当額が 21 万円、老齢厚生年金の基本月額が15 万円（加給年金はない）
\\
※Ｃさんへの影響については、2020 年度（令和2 年度）以降、名目賃金の変動および新規裁
定者の年金額の改定はないものとして解答すること。
\end{itemize}

\end{shortanswer}


\begin{answerbox}
  \begin{enumerate}[(1)]
  \item 65歳以上の在職中の老齢厚生年金受給者の年金額について、従前は被保険者の資格喪失時までは改定されなかったが、
  毎年10月に改定を行うよう変更する。在職定時改定の導入により、就労の継続による年金額の増加を、
  退職をまたずに毎年年金額に反映させることができる。
  \item ＜見直しの内容＞
  
60歳から64歳までに支給される特別支給の老齢厚生年金を対象とした在職老齢年金制度について、年金の支給が停止される基準を28万円から65歳以上の在職老齢年金制度と同じ金額（令和2年度は47万円、令和7年度は51万円）に変更する。

＜Cさんが受ける影響＞

元々は、$(21+15-28)/2=4万円$が支給停止となる。

法改正後は、$21+15<47$となるので、支給停止額は0円となる。

\end{enumerate}
\end{answerbox}

\begin{explanationbox}
  昔の時事問題。
\end{explanationbox}

\begin{fillblank}[年金1 2019 問題1(5)\textcircled{4} 改題]
  \the\year 年度の在職老齢年金の65 歳以上の支給停止調整額は（　　A　　）万円である。
\end{fillblank}


%毎年変える
\begin{answerbox}
  \begin{enumerate}[A: ]
    \item 65
  \end{enumerate}
\end{answerbox}


\begin{explanationbox}
在職老齢年金制度の見直しについて

  \url{https://www.mhlw.go.jp/stf/seisakunitsuite/bunya/0000147284_00022.html}
\end{explanationbox}

\begin{fillblank}[年金1 2019 問題1(5)\textcircled{6}]
  65 歳以上で老齢厚生年金の受給権を有する者（以下、本人）が配偶者の死亡により遺族厚生
年金の受給権を有することになった場合、本人の老齢厚生年金を受給した上で、配偶者の老齢
厚生年金に相当する額×３／４と（　　A　　）のいずれか高い方の年金額と比較して本人の老齢
厚生年金の額が低ければ、その差額が遺族厚生年金として支給される。（ただし、本人は遺族基
礎年金の支給を受けないものとする。）
\end{fillblank}


\begin{answerbox}
  \begin{enumerate}[A: ]
    \item 配偶者の老齢厚生年金に相当する額×１／２＋本人の老齢厚生年金の額×１／２
  \end{enumerate}
\end{answerbox}

\begin{truefalse}[年金1 2019 問題2(2)\textcircled{1}]
以下の文章の正誤を判定し、誤っている場合は正しい内容に改めなさい。\\

同時に2 か所の適用事業所A およびB に使用される第1 号厚生年金被保険者について、同
一の月に適用事業所A から170 万円、適用事業所B から30 万円の賞与が支給された。このと
きの当該被保険者の当該月における標準賞与額は、適用事業所毎に標準賞与額の上限を判断
し、適用事業所A に係る標準賞与額は150 万円、適用事業所B に係る標準賞与額は30 万円と
して決定され、この合計である180 万円となる。
\end{truefalse}


\begin{answerbox}
  \textbf{誤}

  標準賞与額は合算額である200万円から上限を判断し、150万円となる。
\end{answerbox}

\begin{truefalse}[年金1 2019 問題2(2)\textcircled{2}]
以下の文章の正誤を判定し、誤っている場合は正しい内容に改めなさい。\\

  老齢厚生年金の支給繰上げの請求は老齢基礎年金の支給繰上げの請求と同時に行わなければ
ならないが、老齢厚生年金の支給繰下げの申出は必ずしも老齢基礎年金の支給繰下げの申出と
同時に行うことを要しない。
\end{truefalse}

\begin{answerbox}
  \textbf{正}
\end{answerbox}


\begin{truefalse}[年金1 2019 問題2(2)\textcircled{3}]
以下の文章の正誤を判定し、誤っている場合は正しい内容に改めなさい。\\

  老齢厚生年金の加給年金は、原則として、厚生年金保険の被保険者期間が20 年以上ある者
が定額部分支給開始年齢に達したときに、その者によって生計を維持されている65 歳未満の
配偶者または子（18 歳の誕生日の属する年度末を経過していない者または20 歳未満の障害等
級1・2 級に該当する者）があるときに支給される。
\end{truefalse}


\begin{answerbox}
  \textbf{正}
\end{answerbox}

\begin{fillblank}[年金1 2018 問題1(4)\textcircled{4}]
  厚生年金の強制適用となる事業所は、以下の a～c のうち、（　　A　　）。
  \begin{enumerate}[a ]
    \item 常時使用する従業員が３人の個人の事業所
    \item 常時使用する従業員が事業主１人のみである法人の事業所
    \item  国の事業所
  \end{enumerate}

\end{fillblank}


\begin{answerbox}
  b, c
\end{answerbox}

\begin{explanationbox}
  a: 常時5人未満を使用する個人事業所は強制適用事業所ではない。
\end{explanationbox}


\begin{fillblank}[年金1 2018 問題1(4)\textcircled{5}]
老齢厚生年金の繰下げ受給に関し、以下の a～c のうち正しい記述は、（　　A　　）。（なお、
昭和17 年 4 月 2 日以降生まれの方の繰下げ受給について考えるものとする。）
  \begin{enumerate}[a ]
    \item 老齢基礎年金とは別に支給の繰下げ時期を選択できる。
    \item 65 歳以降に被保険者であった場合は、その被保険者であった期間に在職老齢年金制度
    を適用したと仮定した場合に支給される老齢厚生年金の額を繰下げ対象額として、繰下
    げ加算額が計算される。
    \item 65 歳 6 か月の時点では、支給繰下げの申出を行うことができない。
  \end{enumerate}

\end{fillblank}


\begin{answerbox}
  a, b, c
\end{answerbox}

\begin{fillblank}[年金1 2018 問題1(4)\textcircled{6}]
厚生年金の被保険者である夫が死亡したときに（　　A　　）が受ける遺族厚生年金には、夫の
死亡時より中高齢寡婦加算が行われる。（なお、子とは、18 歳到達年度の末日を経過していない
子、または20 歳未満で障害等級 1 級または 2 級の障害の状態にある子をいう。）
\end{fillblank}


\begin{answerbox}
  \begin{enumerate}[A: ]
    \item 40歳以上65歳未満で子のない妻
  \end{enumerate}
\end{answerbox}

\begin{shortanswer}[年金1 2018 問題2(5)\textcircled{2}]
  厚生年金の標準報酬月額の上限改定の考え方については厚生年金保険法に定められているが、その内容
を簡記しなさい。
\end{shortanswer}

\begin{answerbox}
  年度末における全厚生年金被保険者の標準報酬月額の平均額の２倍に相当する額が最高等級の標準報酬月額を上回り、
  その状態が継続すると認められる場合には、政令で、最高等級の上に等級を追加することができる。
\end{answerbox}

\begin{explanationbox}
  ＜参考＞

  令和7年6月13日年金制度改正法により、標準報酬月額の上限改定が行われる。

  \url{https://www.mhlw.go.jp/stf/seisakunitsuite/bunya/0000147284_00024.html}
\end{explanationbox}

\begin{fillblank}[年金1 H28 問題1(2)F]
  老齢厚生年金の支給開始年齢が 65 歳である男性が、70 歳到達後に老齢厚生年金の繰下げ支
給の申出を行った場合における年金額の増額率は（　　A　　）％である。ただし、65 歳から 70 歳ま
での間において繰下げ支給が可能であり、また 65 歳から 70 歳までの間において在職による支
給停止はないものとする。
\end{fillblank}

\begin{answerbox}
  \begin{enumerate}[A: ]
    \item $0.7\% \times 60ヶ月 = 42\%$
  \end{enumerate}
\end{answerbox}

\begin{fillblank}[年金1 H28 問題1(2)G]
  昭和 21 年 4 月 2 日以降に生まれた者の平成 15 年 4 月以降の期間に係る老齢厚生年金報酬比
例部分（厚生年金保険法本則による水準）の給付乗率は 1,000 分の（　　A　　）である。
\end{fillblank}

\begin{answerbox}
  \begin{enumerate}[A: ]
    \item 5.481
  \end{enumerate}
\end{answerbox}

\begin{fillblank}[年金1 H28 問題1(2)I]
  短時間労働者への厚生年金の適用拡大が平成 28 年 10 月より施行されたが、通常
の労働者およびこれに準ずるものを常時（　　A　　）人（\the\year 年10月時点）を超えて使用する事業主の事業所が対象とさ
れた。
\end{fillblank}

%毎年変える
\begin{answerbox}
  \begin{enumerate}[A: ]
    \item 50
  \end{enumerate}
\end{answerbox}

\begin{explanationbox}
  社会保険の加入対象の拡大について

  \url{https://www.mhlw.go.jp/stf/seisakunitsuite/bunya/0000147284_00021.html}
\end{explanationbox}

\begin{shortanswer}[数理人試験 2024 問題3 設問1 改題]
  老齢厚生年金を受給されている者が厚生年金保険の被保険者であるときに、受給されている老齢
厚生年金の基本月額と総報酬月額相当額に応じて年金額の支給停止（在職老齢年金による調整）
が行われる場合がある。

\warekiyear 年度において、在職老齢年金による調整が適用される条件は何か。また、当該調整が適用
される場合、当該適用後の年金月額はどのように計算されるか簡記せよ。
\end{shortanswer}

\begin{answerbox}
  ＜調整が適用される条件＞

基本月額と総報酬月額相当額の合計額が65万円を超過する場合調整が適用される。

＜計算方法＞

$基本月額 - (基本月額 + 総報酬月額相当額 - 65万円) \div 2$
によって計算される。
\end{answerbox}

\begin{shortanswer}[数理人試験 2024 問題3 設問2 改題]
  令和7年6月13日成立の年金制度改正法において短時間労働者に対する厚生年金保険の適用拡大が実施される。ここで言う、
「短時間労働者」の要件がどのように変わるか簡記しなさい。
\end{shortanswer}

\begin{answerbox}
  「短時間労働者」とは従来、(0)従業員51人以上の企業等で働く以下の要件をすべて満たす者のことを指していた。
  \begin{enumerate}[(1)]
    \item 週の所定労働時間が20時間以上
    \item 所定内賃金が月額8.8万円以上
    \item 2ヶ月を超える雇用の見込みがある
    \item 学生ではないこと
  \end{enumerate}

  上記のうち(2)の賃金要件が撤廃される。（いわゆる106万円の壁の撤廃）

  また、(0)の企業規模要件が令和9年10月1日から令和17年10月1日までの間に段階的に撤廃される。
\end{answerbox}


\begin{shortanswer}[数理人試験 2025 問題3 設問3]
  「厚生年金保険法第20条（標準報酬月額）」に定める、最高等級の上に更に等級を加える標準報
酬月額の等級区分の改定を行うことができる場合を簡記せよ。なお、当該改定が適用できる時期に
ついても解答に含めること。
\end{shortanswer}

\begin{answerbox}
  毎年3月31日における全被保険者の標準報酬月額を平均した額の百分の二百に相当する額
が標準報酬月額等級の最高等級の標準報酬月額を超える場合において、その状態が継続
すると認められるときは、その年の9月1日から、標準報酬月額の等級区分を参酌して、
政令で、当該最高等級の上に更に等級を加える標準報酬月額の等級区分の改定を行う
ことができる。
\end{answerbox}

\begin{fillblank}[年金1 2023 問題1(5)ア]
    \noindent \textbf{◯厚生年金保険法}
  \begin{lawstructure}
    \begin{lawarticle}{第四十三条}{年金額}
      老齢厚生年金の額は、被保険者であつた全期間の平均標準報酬額
      （被保険者期間の計算の基礎となる各月の標準報酬月額と標準賞与額に、
      別表各号に掲げる受給権者の区分に応じてそれぞれ当該各号に定める率（以下「再評価率」という。）
      を乗じて得た額の総額を、当該被保険者期間の月数で除して得た額をいう。
      附則第十七条の六第一項及び第二十九条第三項を除き、以下同じ。）の
      千分の五・四八一に相当する額に被保険者期間の月数を乗じて得た額とする。

      \lawparagraph{2}
      受給権者が毎年（　　A　　）（以下この項において「基準日」という。）において被保険者である場合
      （基準日に被保険者の資格を取得した場合を除く。）の老齢厚生年金の額は、
      基準日の（　　B　　）の被保険者であつた期間をその計算の基礎とするものとし、
      基準日の（　　C　　）から、年金の額を改定する。
      ただし、基準日が被保険者の資格を喪失した日から再び被保険者の資格を取得した日までの間に到来し、
      かつ、当該被保険者の資格を喪失した日から再び被保険者の資格を取得した日までの期間が
      一月以内である場合は、基準日の（　　B　　）の被保険者であつた期間を老齢厚生年金の額の
      計算の基礎とするものとし、基準日の（　　C　　）から、年金の額を改定する。
    \end{lawarticle}
  \end{lawstructure}
\end{fillblank}

\begin{answerbox}
  \begin{enumerate}[A: ]
    \item 九月一日
    \item 属する月前
    \item 属する月の翌月
  \end{enumerate}
\end{answerbox}

\begin{fillblank}[年金1 2023 問題1(5)イ(e)]
  厚生年金保険制度の給付の種類に該当するものとし
て選択肢の中から当てはまるものの記号をすべて挙げると（　　A　　）である。

\begin{table}[htbp]
  \centering
  \begin{tabular}{|llll|} \hline
    【選択肢】 & & & \\ 
    （ア）老齢基礎年金 & （イ）老齢厚生年金 & （ウ）付加年金 & （エ）寡婦年金 \\ 
    （オ）老齢手当金 & （カ）障害手当金 & （キ）遺族手当金 & （ク）死亡一時金 \\ \hline
  \end{tabular}
\end{table}
\end{fillblank}

\begin{answerbox}
  （イ）、（カ）
\end{answerbox}

\begin{fillblank}[年金1 2024 問題1(5)ア]
  厚生年金保険では、賞与は労働者が労働の対償として受ける全てのもののうち、 （　　A　　） 月
を超える期間ごとに受けるものとされ、賞与を受けた月における標準賞与額の上限は
（　　B　　） 万円となる。
\end{fillblank}

\begin{answerbox}
  \begin{enumerate}[A: ]
    \item 3
    \item 150
  \end{enumerate}
\end{answerbox}

\begin{fillblank}[年金1 2024 問題1(5)イ]
  2023 年4 月から導入された「特例的な繰下げみなし増額制度」に基づくと、2017 年 4 月 1 日
以降の６５歳に達した日に老齢基礎年金および老齢厚生年金の受給権が発生し、繰下げ受給の
申し出をせず７２歳に達した日に６５歳到達時点の本来の年金（老齢基礎年金および老齢厚生
年金の合計額が220 万円）を遡って請求した場合は、 （　　A　　） 万円を一括で受け取ることが
できる。 （一万円未満を四捨五入して一万円単位で解答すること。 ）
\end{fillblank}

\begin{answerbox}
  \begin{enumerate}[A: ]
    \item 5年前に支給の繰下げを行ったものとみなす。65歳から67歳までの2年間繰り下げた時の増額率は、
    $0.7\% \times 24ヶ月 = 16.8\%$。16.8\%増額した年金を5年分一括で一時金として受給するので、
    $220万円 \times 1.168 \times 5年 = 1284.8万円$。よって、1285万円。
  \end{enumerate}
\end{answerbox}

\begin{fillblank}[年金1 2025 問題1(5)イ 改題]
  令和7 年（2025 年）6 月20 日に公布された「社会経済の変化を踏まえた年金制度の機能強化
のための国民年金法等の一部を改正する等の法律」（令和７年法律第７４号）において、以下の
改正が行われた。
\begin{enumerate}[(1)]
  \item 在職老齢年金制度の見直しにより、支給停止調整額が引き上げられる。\the\year 年度では
（　　A　　） 万円である。
  \item 段階的に標準報酬月額の上限が引き上げられ、最終的に、 （　　B　　） 年9 月1 日から
（　　C　　） 万円が上限となる。
\end{enumerate}

\end{fillblank}

\begin{answerbox}
 \begin{enumerate}[A: ]
  \item 65
  \item 2029
  \item 75
\end{enumerate} 
\end{answerbox}

\begin{shortanswer}[年金1 2025 問題2(3)ア]
  老齢厚生年金のうち、加給年金の支給要件を簡潔に入力しなさい。 
\end{shortanswer}

\begin{answerbox}
  厚生年金保険の被保険者期間が２０年以上ある者が６５歳到達時点に、その者によって生計を
維持されている６５歳未満の配偶者、子（１８歳に達する日以後の最初の３月３１日までの間
にある者、または２０歳未満の障害等級１・２級に該当する者）がいること。
\end{answerbox}

\begin{shortanswer}[年金1 2025 問題2(3)イ]
令和7 年（2025 年）6 月20 日に公布された「社会経済の変化を踏まえた年金制度の機能強化
のための国民年金法等の一部を改正する等の法律」（令和７年法律第７４号）において、遺族厚
生年金の見直しが行われることとなった。配偶者（子はいないものとし、2028 年度末において
４０歳未満であるものとする。）に係る見直し後の支給期間について、簡潔に入力しなさい。な
お、配偶者の性別や死亡時の配偶者の年齢に応じて支給期間が異なる場合には、その区分も含
めて記載すること。また、所得や障害の状態により配慮が必要な方への措置（配慮措置） の内
容については触れなくてよい。
\end{shortanswer}

\begin{answerbox}
  死亡時の配偶者の年齢が、６０歳未満の場合は５年の有期給付、６０歳以上の場合は無期給付
となる。なお、男女で支給期間に相違はない。
\end{answerbox}

\newpage